\cotefile{36}
\scanpage{38}\mspage{- 98 -}\hfill n\textsuperscript{o} 316
% archive sheet number: 37

un automorphisme de $E$.

Si en outre $E$ est une algèbre (non nécessairement associative), et si $u$ est une dérivation (nilpotente) de $E$, alors $e^u$ est un automorphisme de \uline{l'algèbre} $E$. En effet, si $x,y\in E$, on a
\[
u^p(xy)=\sum\nolimits_{r+s=p}\binom{p}{r}\,u^r(x)\,u^s(y)
\]
pour tout entier $p\geqslant 0$ (formule de Leibniz). Il en résulte que :
\[
\begin{aligned}
e^u(xy)&=\sum\nolimits_{p\geqslant 0}\frac{1}{p!}\,u^p(xy)
=\sum\nolimits_{p\geqslant 0}\sum\nolimits_{r+s=p}\frac{u^r(x)}{r!}\,\frac{u^s(y)}{s!}\\
&=\sum\nolimits_{r,s=0}^{\infty}\frac{u^r(x)}{r!}\,\frac{u^s(y)}{s!}
=e^u(x)\,e^u(y)
\end{aligned}
\]
d'où notre assertion.

Soit maintenant $\unsure{\mathbb{g}|g}$ une algèbre de Lie. Si $x$ appartient au radical nilpotent de $\mathbb{g}$, la dérivation $\mathrm{ad}_{\mathbb{g}}\,x$ de $\mathbb{g}$ est nilpotente. On peut donc poser la définition suivante :

\uline{Définition} 5 - \uline{On appelle automorphisme spécial de} $\mathbb{g}$ \uline{un automorphisme de} $\mathbb{g}$ \uline{de la forme} $e^{\mathrm{ad}\,x}$, \uline{où} $x$ \uline{est dans le radical nilpotent de} $\mathbb{g}$.

Il est clair qu'un automorphisme spécial laisse stable tout idéal de $\mathbb{g}$.

\uline{Définition} 6 - \uline{Soient} $\mathbb{g}$ \uline{une algèbre de Lie}, $\unsure{\mathbb{r}|r}$ \uline{son radical}. \uline{On appelle sous-algèbre de Levi de} $\mathbb{g}$ \uline{toute sous-algèbre supplémentaire de} $\mathbb{r}$ \uline{dans} $\mathbb{g}$.

Une sous-algèbre de Levi est isomorphe à $\mathbb{g}/\mathbb{r}$, donc semi-simple. Comme une sous-algèbre semi-simple n'a que $0$ en commun avec $\mathbb{r}$, toute sous-algèbre semi-simple $\unsure{\mathbb{h}|h}$ telle que $\mathbb{g}=\mathbb{r}+\mathbb{h}$ est une sous-algèbre de Levi ; par conséquent, l'image d'une sous-algèbre de Levi par un homomorphisme surjectif est une sous-algèbre de Levi.
